% !TeX root=../main.tex

% Alphabetical page numbering
\cleardoublepage
\pagenumbering{harfi}

% Abstract
\cleardoublepage
\thispagestyle{empty}
\section*{چکیده}
\addcontentsline{toc}{chapter}{چکیده}

انتقال سریع بار در جرثقیل‌های سقفی با چالشی دوگانه روبه‌روست: نوسان‌های نامطلوب آونگ حین حرکت و محدودیت‌های فیزیکی عملگرها. در این پژوهش، به منظور مهار نوسان‌ها از طریق تغییر طول کابل، دینامیک غیرخطی سیستم بر پایه روش اویلر-لاگرانژ استخراج شد و تکینگی کنترل‌پذیری آن در وضعیت تعادل آویزان به اثبات رسید. ارزیابی اولیه عملکرد کنترل‌کننده فیدبک حالت خطی (\lr{LQR}) مشخص کرد که در مانورهای حرکتی بلندبرد (مانند مسیر ۵ متری)، نقض فرض‌های خطی‌سازی و اشباع نیروی محرک به ناپایداری ساختاری و چرخش کامل بار می‌انجامد.

برای مدیریت دقیق این دینامیک، معماری یکپارچه‌ای متشکل از برنامه‌ریز مسیر سینماتیکی نرم (\lr{S-Curve}) و کنترل پیش‌بین مدل غیرخطی (\lr{NMPC}) توسعه یافت. شبیه‌سازی فیزیکی سیستم در محیط \lr{Simscape Multibody} نشان داد که این الگوریتم، قیود سخت شامل سقف نیروی ۵۰ نیوتن، شتاب ۵ متر بر مجذور ثانیه و دامنه نوسان کمتر از ۱۵ درجه را در طول مسیر کاملاً ارضا می‌کند. همچنین، استفاده از عملگر وینچ به عنوان یک مکانیزم میرایی فعال (\lr{Active Damping})، موجب شد تا با تغییر طول کابل حین حرکت، انرژی ارتعاشی مستهلک شده و نوسان‌های پسماند در لحظه توقف کالسکه به طور کامل سرکوب شود.

\vspace{1cm}
\noindent\textbf{واژگان کلیدی:} جرثقیل سقفی، کنترل‌کننده پیش‌بین مدل غیرخطی (\lr{NMPC})، میرایی فعال، طول کابل متغیر، برنامه‌ریز مسیر، تنظیم‌گر خطی درجه دوم (\lr{LQR}).

% Table of Contents
\cleardoublepage
\tableofcontents

% List of Abbreviations
% \cleardoublepage
% \printacronyms

% List of Figures
\cleardoublepage
\listoffigures

% List of Tables
% \cleardoublepage
% \listoftables